\documentclass[10pt,a4paper]{amsart}
\usepackage[margin=19mm]{geometry}
\usepackage{amsmath,amssymb,mathtools}
\usepackage[scr=rsfs,cal=euler]{mathalfa}
\usepackage{xfrac}
\usepackage[unicode,hidelinks,bookmarks=false]{hyperref}
\newcommand{\ie}{i.e.\ }
\newcommand{\cf}{cf.\ }
\newcommand{\resp}{respectively\ }
\newcommand{\Subsection}{\subsection}
\providecommand{\nobreakdash}{\nobreak\hbox{-}}
\setlength{\emergencystretch}{2em}
\multlinegap0pt
\usepackage{bm}
\usepackage{booktabs}
\usepackage{nicefrac}
\usepackage{amssymb}
\usepackage{tcolorbox}
\usepackage{dsfont}
\usepackage{mathtools}
\usepackage[shortlabels]{enumitem}
\usepackage{xcolor}
\usepackage{natbib}
\newtheorem{theorem}{Theorem}
\providecommand{\subsection}{}
\title{Orthogonal parametrisations of Extreme-Value distributions}
\author{Nathan Huet, Ilaria Prosdocimi}
\date{Source: arXiv:2602.16283v1 (CC BY 4.0)}
\begin{document}
\maketitle

\noindent\emph{Source and licence.} Orthogonal parametrisations of Extreme-Value distributions. Version arXiv:2602.16283v1 (CC BY 4.0), distributed by arXiv under the Creative Commons Attribution 4.0 International licence, http://creativecommons.org/licenses/by/4.0/, as recorded in the arXiv licence metadata for this exact version and shown on the abstract page https://arxiv.org/abs/2602.16283v1. Full licence notice: \texttt{licenses/arxiv-2602.16283-CC-BY-4.0.txt}.\par\smallskip

\noindent\textit{Editorial scope.} Counted unit: the article's theorem thm:Gumbel (source0.tex lines 202--216). The article's complete original proof of it is delivered with it (source0.tex lines 590--619, one \texttt{proof} environment of 30 lines, which the article places away from the statement). No mathematics is written, repaired or re-derived: the statement and the proof are the article's own lines, in the article's own notation, and no retained line is substituted or re-flowed.  Retained uncounted as the source-local context the proof needs: lines 170--172; lines 446--456.  Nothing was omitted from any retained span, so the package carries no omission record.  No retained line contains a citation, so the wrapper carries no bibliography and no bibliography entry.  No retained line contains a figure environment, an image-inclusion command or drawing code, so no image is delivered and the package declares no asset dependency.  Editorial formatting only, in addition: an AMS \texttt{amsart} preamble that restates the article's own declarations and macros used by the retained lines, namely the environments \texttt{theorem} on separate counters, together with the standard packages those lines need.  The retained environments print in this document as Theorem 1 in order of appearance, whereas the article prints the same items with the numbering of its own full document.  The complete native source of the article as distributed in the arXiv e-print for this exact version is bundled unchanged as \texttt{sources/194866/source0.tex}.
\begin{equation}\label{eq:orthogonal_pde}
    \sum_{i=1}^p i^*_{\Tilde{\theta}_i\Tilde{\theta}_j} \frac{\partial\Tilde{\theta}_i(\psi, \bm{\lambda})}{\partial\psi} = - i^*_{\psi\Tilde{\theta}_j},
\end{equation}

\begin{equation*}
p(\xi)=(1+\xi)^2 \Gamma(1+2 \xi), \quad q(\xi)=\Gamma(2+\xi)\left(\psi(1+\xi)+\xi^{-1}+1\right),
\end{equation*}
with $\Gamma$ the Gamma function, $\psi(z) = \Gamma'(z)/\Gamma(z)$ the digamma function and $\gamma$ the Euler-Mascheroni constant.

\subsection{Gumbel distribution}\label{sec:gumbel_fisher}

    The Fisher information of the $Gumbel(\mu,\sigma)$ distribution are
\begin{equation*}
    i_{\mu\mu} = 1/\sigma^2; \quad i_{\sigma \sigma} = (\pi^2/6+\gamma^2 - 2 \gamma +1)/\sigma^2; \quad i_{\mu \sigma} = (\gamma - 1)/\sigma^2,
\end{equation*}

\begin{theorem}[Orthogonal reparametrisation of the Gumbel distribution]\label{thm:Gumbel}
    An orthogonal reparametrisation of the Gumbel distribution is given as follows
\begin{enumerate}
    \item with transformed location parameter $\mu$,  $(\nu,\sigma) \mapsto (\Tilde{\mu}(\nu,\sigma),\sigma)$, with
    \begin{equation*}
        \Tilde{\mu}(\nu,\sigma) = (1- \gamma)\sigma + C_1(\nu),
    \end{equation*}
    where $\gamma$ denotes the Euler-Mascheroni constant and where $C_1(\nu)$ is independent of $\sigma$.
    \item with transformed scale parameter $\sigma$, $(\mu,\rho) \mapsto (\mu,\Tilde{\sigma}(\mu,\rho)$, with
    \begin{equation*}
        \Tilde{\sigma}(\mu,\rho) = \frac{1- \gamma}{\pi^2/6+\gamma^2 -2\gamma+1}\mu + C_2(\rho),
    \end{equation*}
    where $C_2(\rho)$ is independent of $\mu$.
\end{enumerate}
\end{theorem}

\begin{proof}[Proof of Theorem~\ref{thm:Gumbel}]
According to Equation~\eqref{eq:orthogonal_pde}, a reparametrisation with a transformed location parameter $(\nu,\sigma) \mapsto (\Tilde{\mu}(\nu,\sigma),\sigma)$, needs to fulfil
\begin{equation*}
    i_{\Tilde{\mu}\Tilde{\mu}} \frac{\partial \Tilde{\mu}(\nu,\sigma)}{\partial \sigma}= - i_{\Tilde{\mu}\sigma}.
\end{equation*}
Given the Fisher information in Appendix~\ref{sec:gumbel_fisher}, this equation rewrites as
\begin{equation*}
    \frac{\partial \Tilde{\mu}(\nu,\sigma)}{\partial \sigma} = 1- \gamma.
\end{equation*}
A solution of this equation is given by
\begin{equation*}
\Tilde{\mu}(\nu,\sigma) = (1- \gamma)\sigma + C_1(\nu),
\end{equation*}
where $C_1(\nu)$ is independent of $\sigma$.


Similarly, a reparametrisation with a transformed scale  $(\mu,\rho) \mapsto (\mu,\Tilde{\sigma}(\mu,\rho))$, needs to fulfil
\begin{equation*}
    i_{\Tilde{\sigma}\Tilde{\sigma}} \frac{\partial \Tilde{\sigma}(\mu,\rho)}{\partial \mu}= - i_{\mu\Tilde{\sigma}}.
\end{equation*}
Given the Fisher information in Appendix~\ref{sec:gumbel_fisher}, this equation rewrites as
\begin{equation*}
    \frac{\partial \Tilde{\sigma}(\mu,\rho)}{\partial \mu} = \frac{1- \gamma}{\pi^2/6+\gamma^2 - 2 \gamma +1}.
\end{equation*}
A solution of this equation is given by
\begin{equation*}
\Tilde{\sigma}(\mu,\rho) = \frac{1- \gamma}{2c_1 -2\gamma+1}\mu + C_2(\rho),
\end{equation*}
where $C_2(\rho)$ is independent of $\mu$.
\end{proof}

\end{document}
